\documentclass[a4paper,12pt]{article}
\usepackage[utf8]{inputenc}
\usepackage[T1]{fontenc}
\usepackage[ngerman]{babel}
\usepackage{amsmath, amssymb}
\usepackage{geometry}
\usepackage{tikz}
\usetikzlibrary{shapes.geometric, arrows.meta, positioning, fit, backgrounds}
\usepackage{parskip}
\usepackage{titlesec}
\usepackage{booktabs}
\usepackage{enumitem}

\geometry{left=3cm, right=3cm, top=2.5cm, bottom=2.5cm}

\titleformat{\section}{\large\bfseries}{}{0em}{}[\titlerule]
\titleformat{\subsection}{\normalsize\bfseries}{}{0em}{}

\tikzset{
  box/.style={rectangle, draw, rounded corners=3pt, text width=4.5cm,
              align=center, minimum height=1.0cm, font=\small, inner sep=4pt},
  boxw/.style={box, text width=5.5cm},
  arrow/.style={-{Latex[length=2mm]}, thick},
}

\begin{document}

\begin{center}
  {\LARGE\bfseries Studierhinweise zu Lektion 5}\\[0.5em]
  {\large Mathematische Grundlagen (Modul 61111)}
\end{center}

\vspace{1em}

Diese Lektion widmet sich dem zentralen Anliegen der reellen Analysis: der
Untersuchung von Funktionen $f: D \to \mathbb{R}$. In Kapitel~14 werden Grundbegriffe
für Funktionen eingeführt (Monotonie, Umkehrfunktion, Beschränktheit) und
wichtige Beispielklassen vorgestellt -- Polynomfunktionen, rationale Funktionen,
allgemeine Potenzfunktionen, Exponential- und Logarithmusfunktion. Kapitel~15
führt den Begriff der \emph{Stetigkeit} ein: Eine stetige Funktion reagiert auf
kleine Änderungen des Arguments nur mit kleinen Änderungen des Funktionswerts.
Sie lernen das $\varepsilon$-$\delta$-Kriterium sowie die Sätze von Bolzano
(Nullstellensatz und Zwischenwertsatz) kennen. Kapitel~16 behandelt die
\emph{Differenzierbarkeit}: Existiert der Grenzwert des Differenzenquotienten,
so nennen wir ihn die Ableitung und interpretieren ihn als Steigung der Tangente.
Die Differentiationsregeln (Produktregel, Quotientenregel, Kettenregel) ermöglichen
das Ableiten aller in der Praxis auftretenden Funktionen.

Lesen Sie den Text aktiv -- rechnen Sie Beweise nach, arbeiten Sie die
Aufgaben selbst durch, und stellen Sie sich bei jeder Definition ein konkretes
Beispiel vor.

% -----------------------------------------------------------------------
\section*{Struktur der Lektion 5}

\begin{center}
\begin{tikzpicture}[node distance=0.5cm and 0.7cm]

  % Column headers
  \node[font=\bfseries\small] (h1) at (0,0)      {Kap.\ 14: Funktionen};
  \node[font=\bfseries\small] (h2) at (5.5,0)    {Kap.\ 15: Stetigkeit};
  \node[font=\bfseries\small] (h3) at (11,0)     {Kap.\ 16: Differenzierbarkeit};

  % Kapitel 14
  \node[box, below=0.8cm of h1] (f1) {14.1 Grundbegriffe\\Graph, Umkehrfunktion,\\Monotonie, Beschränktheit};
  \node[box, below=0.5cm of f1] (f2) {14.2 Beispiele:\\Polynom- und rationale\\Funktionen};
  \node[box, below=0.5cm of f2] (f3) {14.2 Allg.\ Potenz $x^\rho$\\Exponentialfunktion $a^x$};
  \node[box, below=0.5cm of f3] (f4) {14.2 Logarithmus $\log_a$\\natürl.\ Logarithmus $\ln$};

  % Kapitel 15
  \node[box, below=0.8cm of h2] (s1) {15.1 Stetigkeit in $a$\\$\varepsilon$-$\delta$-Kriterium\\Rechenregeln};
  \node[box, below=0.5cm of s1] (s2) {15.2 Stetigkeit auf\\Intervallen: Bolzano,\\Zwischenwertsatz, Min/Max};
  \node[box, below=0.5cm of s2] (s3) {15.3 Grenzwerte von\\Funktionen $\lim_{x\to a}f(x)$\\hebbare Unstetigkeit};

  % Kapitel 16
  \node[box, below=0.8cm of h3] (d1) {16.1 Ableitung $f'(a)$\\Differenzenquotient\\lokale Extrema};
  \node[box, below=0.5cm of d1] (d2) {16.2 Differentiationsregeln\\Produkt-, Quotienten-,\\Kettenregel, Umkehrregel};
  \node[box, below=0.5cm of d2] (d3) {16.3 Ableitungen:\\Polynom, exp, ln, $x^a$\\Sinus, Cosinus};

  % Pfeile Kap 14
  \draw[arrow] (f1) -- (f2);
  \draw[arrow] (f2) -- (f3);
  \draw[arrow] (f3) -- (f4);

  % Pfeile Kap 15
  \draw[arrow] (s1) -- (s2);
  \draw[arrow] (s2) -- (s3);

  % Pfeile Kap 16
  \draw[arrow] (d1) -- (d2);
  \draw[arrow] (d2) -- (d3);

  % Kapitelübergreifende Pfeile
  \draw[arrow] (f4.east) -- ++(0.3,0) |- (s1.west);
  \draw[arrow] (s3.east) -- ++(0.3,0) |- (d1.west);

\end{tikzpicture}
\end{center}

\newpage

% -----------------------------------------------------------------------
\section*{Zielelement 14.1 -- Grundbegriffe für Funktionen}

\subsection*{Lerninhalte}

\begin{center}
\begin{tikzpicture}[node distance=0.5cm and 1.0cm]
  \node[box] (def) {Funktion $f: D \to \mathbb{R}$, Graph, Bild};
  \node[box, below=0.5cm of def] (umk) {Umkehrfunktion $f^{-1}$\\(existiert $\Leftrightarrow$ $f$ injektiv)};
  \node[box, below=0.5cm of umk] (mon) {Monotonie (wachsend/fallend)\\streng monoton};
  \node[box, below=0.5cm of mon] (bes) {Beschränktheit von Funktionen};
  \node[box, below=0.5cm of bes] (res) {Restriktion $f|_{D'}$};
  \draw[arrow] (def) -- (umk);
  \draw[arrow] (umk) -- (mon);
  \draw[arrow] (mon) -- (bes);
  \draw[arrow] (bes) -- (res);
\end{tikzpicture}
\end{center}

\subsection*{Lernziele}

Nach Durcharbeiten dieses Abschnitts sollten Sie:
\begin{itemize}
  \item den Unterschied zwischen dem Analysis- und dem Lineare-Algebra-Begriff
        der Komposition kennen (Bild von $f$ muss Teilmenge des Definitionsbereichs
        von $g$ sein, nicht gleich),
  \item erklären können, wann eine Umkehrfunktion existiert, und den Graphen von
        $f^{-1}$ durch Spiegelung an der Diagonale $d = \{(z,z)\}$ konstruieren,
  \item die Begriffe monoton wachsend, streng monoton fallend usw.\ an Beispielen
        überprüfen und Proposition~14.1.11 (streng monotone Funktionen sind injektiv)
        anwenden können,
  \item beschränkte Funktionen erkennen und algebraische Verknüpfungen beschränkter
        Funktionen auf Beschränktheit untersuchen.
\end{itemize}

\subsection*{Selbstkontrollelement 14.1}

Sei $f: \mathbb{R} \to \mathbb{R}$, $f(x) = 2x + 1$. Bestimmen Sie die Umkehrfunktion
$f^{-1}$ und beschreiben Sie, wie man den Graphen von $f^{-1}$ aus dem Graphen von
$f$ erhält.

\newpage

% -----------------------------------------------------------------------
\section*{Zielelement 14.2a -- Polynomfunktionen und rationale Funktionen}

\subsection*{Lerninhalte}

\begin{center}
\begin{tikzpicture}[node distance=0.5cm and 1.0cm]
  \node[box] (pol) {Polynom $p \in K[T]$:\\Produkt, Division mit Rest};
  \node[box, below=0.5cm of pol] (nst) {Nullstellen und Faktor $(T-\lambda)$\\Anzahl der Nullstellen $\leq \deg(p)$};
  \node[box, below=0.5cm of nst] (pf)  {Polynomfunktion $\tilde{p}: \mathbb{R} \to \mathbb{R}$\\Identitätssatz};
  \node[box, below=0.5cm of pf]  (rat) {Rationale Funktion $\tilde{p}/\tilde{q}$\\Kürzen, Definitionsbereich};
  \draw[arrow] (pol) -- (nst);
  \draw[arrow] (nst) -- (pf);
  \draw[arrow] (pf) -- (rat);
\end{tikzpicture}
\end{center}

\subsection*{Lernziele}

Nach Durcharbeiten dieses Abschnitts sollten Sie:
\begin{itemize}
  \item Polynome in $K[T]$ multiplizieren und mit Rest dividieren können,
  \item Proposition~14.2.8 (Nullstellen $\Leftrightarrow$ Teiler $(T-\lambda)$) und
        Korollar~14.2.9 (höchstens $n$ Nullstellen) kennen und anwenden,
  \item den Identitätssatz (Satz~14.2.13) formulieren: Stimmen zwei Polynomfunktionen
        in mehr als $\max(\deg p, \deg q)$ Punkten überein, so sind die Polynome gleich,
  \item rationale Funktionen kürzen und ihren Definitionsbereich bestimmen.
\end{itemize}

\subsection*{Selbstkontrollelement 14.2a}

Zeigen Sie, dass $\lambda = 2$ eine Nullstelle von $p = T^3 - 6T + 4$ ist, und
geben Sie die Division $p = (T - 2) \cdot q + r$ explizit an.

\newpage

% -----------------------------------------------------------------------
\section*{Zielelement 14.2b -- Allgemeine Potenz und Exponentialfunktion}

\subsection*{Lerninhalte}

\begin{center}
\begin{tikzpicture}[node distance=0.5cm and 1.0cm]
  \node[box] (rat) {Rationale Potenzen $a^r$, $r \in \mathbb{Q}$};
  \node[box, below=0.5cm of rat] (rea) {Allgemeine Potenz $a^\rho$, $\rho \in \mathbb{R}$\\als Grenzwert $(a^{r_n})$ mit $r_n \to \rho$};
  \node[box, below=0.5cm of rea] (prg) {Potenzregeln für $\rho, \sigma \in \mathbb{R}$\\$(a^\rho)^\sigma = a^{\rho\sigma}$};
  \node[box, below=0.5cm of prg] (pot) {Potenzfunktion $x \mapsto x^\rho$: Monotonie, Bild, Umkehrung};
  \node[box, below=0.5cm of pot] (exp) {Exponentialfunktion $\mathrm{exp}_a(x) = a^x$\\Eigenschaften, Graph};
  \node[box, below=0.5cm of exp] (eu)  {Eulersche Zahl $e$, $\mathrm{exp} = \mathrm{exp}_e$\\$\lim (1+x_n)^{1/x_n} = e$};
  \draw[arrow] (rat) -- (rea);
  \draw[arrow] (rea) -- (prg);
  \draw[arrow] (prg) -- (pot);
  \draw[arrow] (pot) -- (exp);
  \draw[arrow] (exp) -- (eu);
\end{tikzpicture}
\end{center}

\subsection*{Lernziele}

Nach Durcharbeiten dieses Abschnitts sollten Sie:
\begin{itemize}
  \item erklären können, wie $a^\rho$ für irrationales $\rho$ als Grenzwert
        einer Folge mit rationalen Exponenten definiert wird,
  \item die fünf Potenzregeln (Satz~14.2.32) auswendig kennen und anwenden,
  \item die Monotonieeigenschaften der allgemeinen Potenzfunktion (Proposition~14.2.25)
        und der Exponentialfunktion (Satz~14.2.35) kennen,
  \item den Grenzwert $\lim_{n\to\infty}(1+x_n)^{1/x_n} = e$ (Satz~14.2.39) anwenden.
\end{itemize}

\subsection*{Selbstkontrollelement 14.2b}

Berechnen Sie $2^\pi$ schrittweise: Geben Sie eine Folge rationaler Zahlen $(r_n)$
mit $r_n \to \pi$ an und erläutern Sie, warum $(2^{r_n})$ konvergiert.

\newpage

% -----------------------------------------------------------------------
\section*{Zielelement 14.2c -- Logarithmus}

\subsection*{Lerninhalte}

\begin{center}
\begin{tikzpicture}[node distance=0.5cm and 1.0cm]
  \node[box] (log) {Logarithmus $\log_a$:\\Umkehrfunktion von $\mathrm{exp}_a$};
  \node[box, below=0.5cm of log] (eig) {Eigenschaften: Monotonie, Bild $= \mathbb{R}$\\$\log_a(1) = 0$, $\log_a(a) = 1$};
  \node[box, below=0.5cm of eig] (rec) {Rechenregeln:\\$\log_a(xy) = \log_a x + \log_a y$\\$\log_a(x^\rho) = \rho\log_a x$};
  \node[box, below=0.5cm of rec] (ln)  {Natürlicher Logarithmus $\ln = \log_e$\\$x^a = \exp(a\ln x)$,\; $\exp_a(x) = \exp(x\ln a)$};
  \draw[arrow] (log) -- (eig);
  \draw[arrow] (eig) -- (rec);
  \draw[arrow] (rec) -- (ln);
\end{tikzpicture}
\end{center}

\subsection*{Lernziele}

Nach Durcharbeiten dieses Abschnitts sollten Sie:
\begin{itemize}
  \item den Logarithmus als Umkehrfunktion der Exponentialfunktion definieren können,
  \item die Eigenschaften des Logarithmus (Satz~14.2.42) kennen und einsetzen können,
  \item die Rechenregeln für den Logarithmus sicher anwenden (insbesondere die
        Logarithmusgesetze für Produkte, Quotienten und Potenzen),
  \item $\log_a(x)$, $\exp_a(x)$ und die allgemeine Potenzfunktion durch $\exp$ und
        $\ln$ ausdrücken können (Proposition~14.2.44).
\end{itemize}

\subsection*{Selbstkontrollelement 14.2c}

Berechnen Sie $\log_2(8)$ und vereinfachen Sie $\log_a(a^3 \cdot \sqrt{a})$ so
weit wie möglich.

\newpage

% -----------------------------------------------------------------------
\section*{Zielelement 15.1 -- Stetigkeit: Grundlagen}

\subsection*{Lerninhalte}

\begin{center}
\begin{tikzpicture}[node distance=0.5cm and 1.0cm]
  \node[box] (def) {Stetigkeit in $a$: $(a_n) \to a$ $\Rightarrow$ $(f(a_n)) \to f(a)$};
  \node[box, below=0.5cm of def] (bsp) {Beispiele: Dirichletfunktion unstetig,\\$x^\rho$, $\exp_a$, $\log_a$ stetig};
  \node[box, below=0.5cm of bsp] (eed) {$\varepsilon$-$\delta$-Kriterium für Stetigkeit\\(Satz 15.1.9)};
  \node[box, below=0.5cm of eed] (rec) {Rechenregeln: $f \pm g$, $fg$, $\frac{f}{g}$, $h \circ f$ stetig};
  \node[box, below=0.5cm of rec] (pol) {Polynomfunktionen und rationale\\Funktionen sind stetig};
  \draw[arrow] (def) -- (bsp);
  \draw[arrow] (bsp) -- (eed);
  \draw[arrow] (eed) -- (rec);
  \draw[arrow] (rec) -- (pol);
\end{tikzpicture}
\end{center}

\subsection*{Lernziele}

Nach Durcharbeiten dieses Abschnitts sollten Sie:
\begin{itemize}
  \item die Folgendefinition der Stetigkeit (Definition~15.1.1) und das
        $\varepsilon$-$\delta$-Kriterium (Satz~15.1.9) formulieren und ineinander
        überführen können,
  \item die Äquivalenz beider Formulierungen beweisen können,
  \item die Rechenregeln für stetige Funktionen (Proposition~15.1.11) anwenden,
  \item konkrete Funktionen auf Stetigkeit oder Unstetigkeit prüfen (inkl.\
        Dirichletfunktion und Größte-Ganze-Funktion).
\end{itemize}

\subsection*{Selbstkontrollelement 15.1}

Beweisen Sie mit dem $\varepsilon$-$\delta$-Kriterium, dass
$f: \mathbb{R} \to \mathbb{R}$, $f(x) = 3x - 1$, stetig ist.

\newpage

% -----------------------------------------------------------------------
\section*{Zielelement 15.2 -- Stetige Funktionen auf Intervallen}

\subsection*{Lerninhalte}

\begin{center}
\begin{tikzpicture}[node distance=0.5cm and 1.0cm]
  \node[box] (bol) {Nullstellensatz von Bolzano\\$f(a) < 0 < f(b)$ $\Rightarrow$ Nullstelle};
  \node[box, below=0.5cm of bol] (zwi) {Zwischenwertsatz:\\$f$ nimmt jeden Wert zw.\ $f(a)$ und $f(b)$ an};
  \node[box, below=0.5cm of zwi] (mm)  {Satz vom Minimum und Maximum:\\$f: [a,b] \to \mathbb{R}$ stetig $\Rightarrow$ Extrema};
  \node[box, below=0.5cm of mm]  (mon) {injektiv und stetig auf Intervall\\$\Leftrightarrow$ streng monoton};
  \node[box, below=0.5cm of mon] (inv) {Umkehrfunktionen stetiger Funktionen\\sind stetig};
  \draw[arrow] (bol) -- (zwi);
  \draw[arrow] (zwi) -- (mm);
  \draw[arrow] (mm) -- (mon);
  \draw[arrow] (mon) -- (inv);
\end{tikzpicture}
\end{center}

\subsection*{Lernziele}

Nach Durcharbeiten dieses Abschnitts sollten Sie:
\begin{itemize}
  \item den Nullstellensatz und den Zwischenwertsatz von Bolzano formulieren und
        beweisen können,
  \item den Satz vom Minimum und Maximum (Korollar~15.2.9) kennen und anwenden,
  \item das Lemma~15.2.10 (stetig und injektiv auf $[a,b]$ $\Rightarrow$ streng monoton)
        kennen sowie dessen Beweis und Konsequenzen verstehen,
  \item erklären, warum die Umkehrfunktion einer auf einem Intervall stetigen und
        injektiven Funktion ebenfalls stetig ist (Proposition~15.2.14).
\end{itemize}

\subsection*{Selbstkontrollelement 15.2}

Sei $f: [0,1] \to \mathbb{R}$ stetig mit $f(0) = -1$ und $f(1) = 3$. Welche Aussage
liefert der Zwischenwertsatz? Welche Werte nimmt $f$ garantiert an?

\newpage

% -----------------------------------------------------------------------
\section*{Zielelement 15.3 -- Grenzwerte von Funktionen}

\subsection*{Lerninhalte}

\begin{center}
\begin{tikzpicture}[node distance=0.5cm and 1.0cm]
  \node[box] (hp)  {Häufungspunkt von $D \subseteq \mathbb{R}$};
  \node[box, below=0.5cm of hp] (kon) {Konvergenz von $f$ in $a$:\\$(f(a_n))$ konvergiert für alle $(a_n) \to a$};
  \node[box, below=0.5cm of kon] (grw) {Grenzwert $\lim_{x \to a} f(x)$\\(unabhängig von der Folge)};
  \node[box, below=0.5cm of grw] (ste) {$f$ stetig in $a$ $\Leftrightarrow$ $\lim_{x\to a}f(x) = f(a)$};
  \node[box, below=0.5cm of ste] (heb) {Hebbare Unstetigkeit\\$\varepsilon$-$\delta$-Kriterium für Grenzwert};
  \draw[arrow] (hp) -- (kon);
  \draw[arrow] (kon) -- (grw);
  \draw[arrow] (grw) -- (ste);
  \draw[arrow] (ste) -- (heb);
\end{tikzpicture}
\end{center}

\subsection*{Lernziele}

Nach Durcharbeiten dieses Abschnitts sollten Sie:
\begin{itemize}
  \item den Begriff des Häufungspunkts und die Definition der Konvergenz einer
        Funktion in einem Punkt (Definition~15.3.4) erklären können,
  \item den Zusammenhang zwischen Stetigkeit und dem Grenzwert $\lim_{x\to a}f(x)$
        (Proposition~15.3.9) kennen,
  \item das $\varepsilon$-$\delta$-Kriterium für Funktionengrenzwerte (Satz~15.3.15)
        anwenden und hebbare Unstetigkeiten erkennen und beheben können,
  \item die Rechenregeln für Funktionenkonvergenz (Proposition~15.3.17) anwenden.
\end{itemize}

\subsection*{Selbstkontrollelement 15.3}

Berechnen Sie $\displaystyle\lim_{x \to 2} \frac{x^2 - 4}{x - 2}$ und erklären
Sie, ob die Funktion $f(x) = \frac{x^2-4}{x-2}$ eine hebbare Unstetigkeit in
$x = 2$ besitzt.

\newpage

% -----------------------------------------------------------------------
\section*{Zielelement 16.1 -- Die Ableitung einer Funktion}

\subsection*{Lerninhalte}

\begin{center}
\begin{tikzpicture}[node distance=0.5cm and 1.0cm]
  \node[box] (dif) {Differenzierbarkeit in $a$:\\$\lim_{x \to a} \frac{f(x)-f(a)}{x-a} = f'(a)$};
  \node[box, below=0.5cm of dif] (geo) {Geometrische Interpretation:\\Steigung der Tangente, Sekante};
  \node[box, below=0.5cm of geo] (dif2s) {differenzierbar $\Rightarrow$ stetig};
  \node[box, below=0.5cm of dif2s] (ext) {Lokale und globale Extrema\\Definition innerer Punkt};
  \node[box, below=0.5cm of ext] (kri) {Notwendige Bedingung: $f'(a) = 0$\\an inneren lokalen Extremstellen};
  \draw[arrow] (dif) -- (geo);
  \draw[arrow] (geo) -- (dif2s);
  \draw[arrow] (dif2s) -- (ext);
  \draw[arrow] (ext) -- (kri);
\end{tikzpicture}
\end{center}

\subsection*{Lernziele}

Nach Durcharbeiten dieses Abschnitts sollten Sie:
\begin{itemize}
  \item die Definition der Differenzierbarkeit (Definition~16.1.1) durch den
        Grenzwert des Differenzenquotienten formulieren und den geometrischen
        Zusammenhang mit der Tangente erklären können,
  \item beweisen können, dass Differenzierbarkeit Stetigkeit impliziert
        (Proposition~16.1.3),
  \item lokale und globale Extrema definieren und die notwendige Bedingung
        $f'(a) = 0$ an inneren Extremstellen (Proposition~16.1.12) kennen und
        anwenden können.
\end{itemize}

\subsection*{Selbstkontrollelement 16.1}

Beweisen Sie direkt aus der Definition, dass $f(x) = x^2$ in $a = 3$ differenzierbar
ist und berechnen Sie $f'(3)$.

\newpage

% -----------------------------------------------------------------------
\section*{Zielelement 16.2 -- Differentiationsregeln}

\subsection*{Lerninhalte}

\begin{center}
\begin{tikzpicture}[node distance=0.5cm and 1.0cm]
  \node[box] (lin) {Summen- und Faktorregel\\$(f \pm g)' = f' \pm g'$, $(\alpha f)' = \alpha f'$};
  \node[box, below=0.5cm of lin] (pro) {Produktregel\\$(fg)' = f'g + fg'$};
  \node[box, below=0.5cm of pro] (quo) {Quotientenregel und Reziprokregel\\$\left(\frac{f}{g}\right)' = \frac{f'g - fg'}{g^2}$};
  \node[box, below=0.5cm of quo] (ket) {Kettenregel\\$(f \circ g)'(a) = f'(g(a)) \cdot g'(a)$};
  \node[box, below=0.5cm of ket] (umk) {Ableitung der Umkehrfunktion\\$(f^{-1})'(b) = \frac{1}{f'(f^{-1}(b))}$};
  \draw[arrow] (lin) -- (pro);
  \draw[arrow] (pro) -- (quo);
  \draw[arrow] (quo) -- (ket);
  \draw[arrow] (ket) -- (umk);
\end{tikzpicture}
\end{center}

\subsection*{Lernziele}

Nach Durcharbeiten dieses Abschnitts sollten Sie:
\begin{itemize}
  \item alle Differentiationsregeln (Proposition~16.2.1) auswendig kennen und
        beherrschen -- Summe, Produkt, Quotient, Faktor,
  \item die Kettenregel (Satz~16.2.3) sicher anwenden können, auch bei mehrfach
        verschachtelten Kompositionen,
  \item die Formel für die Ableitung der Umkehrfunktion (Proposition~16.2.4) kennen
        und geometrisch interpretieren können.
\end{itemize}

\subsection*{Selbstkontrollelement 16.2}

Leiten Sie $f(x) = \dfrac{x^2 + 1}{x - 1}$ mit der Quotientenregel ab und
$g(x) = (x^3 + 1)^5$ mit der Kettenregel.

\newpage

% -----------------------------------------------------------------------
\section*{Zielelement 16.3 -- Ableitungen wichtiger Funktionen}

\subsection*{Lerninhalte}

\begin{center}
\begin{tikzpicture}[node distance=0.5cm and 1.0cm]
  \node[box] (pw) {Potenzregel $(x^k)' = kx^{k-1}$, $k \in \mathbb{Z}$};
  \node[box, below=0.5cm of pw] (ex) {Exponentialfunktion:\\$(\exp_a(x))' = \exp_a(x) \cdot \ln(a)$,\; $(e^x)' = e^x$};
  \node[box, below=0.5cm of ex] (ln) {Logarithmus:\\$\ln'(x) = \frac{1}{x}$,\; $\log_a'(x) = \frac{1}{x\ln(a)}$};
  \node[box, below=0.5cm of ln] (al) {Allgemeine Potenzfunktion:\\$(x^a)' = ax^{a-1}$};
  \node[box, below=0.5cm of al] (sc) {Winkelfunktionen:\\$\sin' = \cos$,\; $\cos' = -\sin$};
  \draw[arrow] (pw) -- (ex);
  \draw[arrow] (ex) -- (ln);
  \draw[arrow] (ln) -- (al);
  \draw[arrow] (al) -- (sc);
\end{tikzpicture}
\end{center}

\subsection*{Lernziele}

Nach Durcharbeiten dieses Abschnitts sollten Sie:
\begin{itemize}
  \item alle Standard-Ableitungsregeln auswendig wissen: Potenzregel für
        $k \in \mathbb{Z}$ und $a \in \mathbb{R}$, Exponentialfunktion,
        natürlicher Logarithmus, Sinus und Cosinus,
  \item Ableitungen zusammengesetzter Funktionen durch Kombination der
        Differentiationsregeln mit diesen Grundableitungen berechnen,
  \item die Eigenschaften von Sinus und Cosinus (SinCos-1 bis SinCos-5)
        kennen und aus ihnen die Differenzierbarkeit herleiten können.
\end{itemize}

\subsection*{Selbstkontrollelement 16.3}

Berechnen Sie die Ableitungen von
$f(x) = x^4 \exp(x)$,\quad $g(x) = \ln(\sin(x))$ und \quad $h(x) = x^{\sqrt{2}}$.

\end{document}
